% !TeX program = xelatex
% Προετοιμασία υπεράσπισης, σε μορφή ερωτήσεων--απαντήσεων.
% Μεταγλώττιση: xelatex defense_qa.tex  x2  (από μέσα από docs/)
\documentclass[11pt,a4paper]{article}

\usepackage{fontspec}
\usepackage{polyglossia}
\setmainlanguage{greek}
\setotherlanguage{english}
\setmainfont{Linux Libertine O}
\setsansfont{Linux Biolinum O}
\setmonofont{DejaVu Sans Mono}[Scale=0.85]

\usepackage[a4paper,margin=2.4cm]{geometry}
\usepackage{parskip}
\usepackage{amsmath,amssymb}
\usepackage{enumitem}
\usepackage{xcolor}
\definecolor{qcol}{HTML}{1a4f8f}
\definecolor{warn}{HTML}{8f3a1a}
\usepackage[colorlinks=true,urlcolor=qcol,linkcolor=qcol]{hyperref}
\setlength{\emergencystretch}{3em}
\newcommand{\en}[1]{\textenglish{#1}}

% Ερώτηση / απάντηση
\newcommand{\qa}[2]{%
  \par\medskip\noindent{\color{qcol}\textbf{Ε.}}\ \textbf{#1}\par
  \noindent{\color{qcol}\textbf{Α.}}\ #2\par}
% Παγίδα: ερώτηση που θέλει προσοχή
\newcommand{\trap}[1]{\par\smallskip\noindent{\color{warn}\textbf{Προσοχή:}} \emph{#1}\par}

\title{\textbf{Τι πρέπει να ξέρω για να τα εξηγήσω}\\[6pt]
\large\normalfont Προετοιμασία σε μορφή ερωτήσεων -- απαντήσεων}
\author{}
\date{\today}

\begin{document}
\maketitle

\noindent\small
Οι απαντήσεις είναι γραμμένες για να ειπωθούν, όχι για να διαβαστούν: σύντομες, με το
κρίσιμο νούμερο μέσα. Όπου μια ερώτηση είναι παγίδα ή απαιτεί παραδοχή αδυναμίας, αυτό
σημειώνεται ρητά --- η χειρότερη απάντηση σε τέτοια ερώτηση είναι η άμυνα.
\normalsize

% =====================================================================
\section{Το πλαίσιο: γιατί έτσι και όχι αλλιώς}
% =====================================================================

\qa{Γιατί αντιμετωπίζεις τους bookmakers ως experts αντί να φτιάξεις δικό σου μοντέλο
πρόβλεψης;}
{Γιατί δεν θα κέρδιζα. Οι εταιρείες έχουν περισσότερα δεδομένα, περισσότερους πόρους και
οικονομικό κίνητρο να προβλέπουν σωστά. Η τιμή τους ενσωματώνει ήδη πληροφορία που δεν
υπάρχει σε κανένα δημόσιο στατιστικό, όπως τραυματισμοί και ενδεκάδες. Αντί να ανταγωνιστώ
αυτή την πληροφορία, τη δέχομαι ως είσοδο και ρωτάω κάτι διαφορετικό: \emph{ποιον να
πιστέψω και πόσο}. Αυτό είναι γνήσιο πρόβλημα, γιατί οι εταιρείες διαφέρουν συστηματικά σε
ποιότητα.}

\qa{Γιατί online μάθηση και όχι κανονικό train/test;}
{Τρεις λόγοι. Πρώτον, κάθε πρόβλεψη είναι εξ ορισμού εκτός δείγματος, αφού χρησιμοποιεί μόνο
τους προηγούμενους γύρους --- δεν χρειάζεται διαχωρισμός για να είναι τίμια. Δεύτερον, το
regret δεν προϋποθέτει ούτε σταθερή κατανομή ούτε ανεξάρτητα δείγματα, και εδώ καμία από τις
δύο δεν ισχύει: η αγορά αλλάζει μέσα στη δεκαετία και διαδοχικοί αγώνες μοιράζονται
αγωνιστική και πρωτάθλημα. Τρίτον, είναι ο ρεαλιστικός τρόπος λειτουργίας --- οι αποδόσεις
έρχονται μία μία.}

\qa{Τι ακριβώς εγγυάται το regret, και τι δεν εγγυάται;}
{Εγγυάται ότι δεν θα τα πάω ουσιωδώς χειρότερα από τον \emph{καλύτερο σταθερό συνδυασμό} των
experts, για οποιαδήποτε ακολουθία, χωρίς καμία στατιστική υπόθεση. Δεν εγγυάται ότι θα
προβλέψω καλά. Αν όλοι οι experts είναι κακοί, η εγγύηση τηρείται και το αποτέλεσμα είναι
κακό. Είναι \emph{σχετική} εγγύηση, και αυτό πρέπει να το πω πριν μου το πουν.}

\qa{Γιατί log-loss και όχι ακρίβεια;}
{Γιατί η ακρίβεια δεν είναι κατάλληλος κανόνας βαθμολόγησης: αγνοεί το πόσο βέβαιη ήταν η
πρόβλεψη. Και εμπειρικά το έδειξα, δεν το υπέθεσα: στην ίδια μπαταρία ελέγχων το Brier
συμφωνεί με το log-loss σε 18 από 18 ετυμηγορίες, ενώ η ακρίβεια δίνει \textbf{0 σημαντικά
αποτελέσματα στα 18}. Με ακρίβεια γύρω στο 50,6\% για όλες τις μεθόδους, μια μετρική που
πετά τη βεβαιότητα πετά ακριβώς το σήμα που τις ξεχωρίζει.}

% =====================================================================
\section{Οι αλγόριθμοι}
% =====================================================================

\qa{Πες μου σε μία πρόταση τον καθένα από τους τρεις.}
{Ο \textbf{Hedge} κρατά ένα βάρος ανά expert και το μειώνει εκθετικά ανάλογα με την απώλειά
του. Ο \textbf{OGD} κάνει ένα βήμα αντίθετα στην κλίση και προβάλλει πίσω στο simplex. Ο
\textbf{FTRL} ξαναϋπολογίζει κάθε φορά τα βάρη από ολόκληρη τη σωρευμένη ιστορία συν έναν
όρο κανονικοποίησης --- «τεμπέλης» αντί «άπληστος».}

\qa{Γιατί τρεις και όχι ένας;}
{Γιατί καλύπτουν τις δύο βασικές οικογένειες OCO, πολλαπλασιαστικά βάρη και βήμα κλίσης, και
επειδή η σύγκρισή τους παράγει ευρήματα. Το σημαντικότερο: ο OGD είναι ο μόνος που δεν χάνει
από τον απλό μέσο όρο σε ομοιογενές panel, και αυτό είναι το ισχυρότερο επιχείρημα υπέρ του.}

\qa{Τι είναι η αναγωγή Sleeping Experts και γιατί τη χρειάζεσαι;}
{Οι κλασικοί αλγόριθμοι υποθέτουν ότι κάθε expert δίνει πρόβλεψη κάθε γύρο. Εδώ η κάλυψη
πηγαίνει από 7\% ως 100\%. Η αναγωγή λέει: κάθε expert κρατά μόνιμη κατάσταση, η πρόβλεψη
σχηματίζεται μόνο από τους ενεργούς με επανακανονικοποίηση, και όσοι κοιμούνται παγώνουν.
Το κέρδος είναι ότι η εγγύηση regret κάθε expert μετριέται \emph{μόνο στους γύρους όπου ήταν
ενεργός}, οπότε ένας bookmaker με 7\% κάλυψη δεν αραιώνει την εγγύηση των υπολοίπων.}

\qa{Γιατί ο OGD θέλει βήμα 100 φορές μεγαλύτερο ενώ ο Hedge 4 φορές μικρότερο;}
{Είναι δομικό, όχι τυχαίο. Στον OGD το $\eta$ \emph{πολλαπλασιάζει} την κλίση ενός γύρου,
οπότε μεγαλύτερο σημαίνει απλώς μεγαλύτερο βήμα. Στον Hedge το $\eta$ βρίσκεται σε
\emph{εκθέτη} πάνω στην απώλεια ενός γύρου, οπότε μικρότερο σημαίνει βραδύτερη απόσβεση,
δηλαδή \emph{μακρύτερη μνήμη}. Παρότι λέγονται και τα δύο «ρυθμός μάθησης», δεν ελέγχουν το
ίδιο μέγεθος.}

\qa{Από πού ξέρεις ότι το $\times100$ δεν είναι overfitting;}
{Τρία πράγματα. Είναι \emph{εσωτερικό} βέλτιστο, όχι άκρο του πλέγματος: το 200, το 500 και
το 1000 είναι χειρότερα. Το train και η ουρά ελέγχου συμφωνούν και τα δύο στο 100. Και
υπάρχει ανεξάρτητη θεωρητική εξήγηση για μέρος του: το φράγμα $M=\log(1/\varepsilon)=13{,}82$
υποθέτει ότι κάποιος μπορεί να τιμολογήσει το πραγματικό αποτέλεσμα στο $10^{-6}$, ενώ στην
πράξη η μέγιστη παρατηρούμενη απώλεια είναι 3,76.}

% =====================================================================
\section{Οι δύσκολες ερωτήσεις}
% =====================================================================

\qa{Το μείγμα δεν αναπαράγει απλώς την Pinnacle, που είναι ούτως ή άλλως μέσα του;}
{Το έλεγξα αφαιρώντας και τις δύο στήλες της και επανεκπαιδεύοντας από την αρχή. Στο πλήρες
panel όλες οι μέθοδοι παραμένουν πάνω από όλους τους εναπομείναντες bookmakers και ο OGD
κερδίζει σημαντικά και τους πέντε. Άρα όχι. \emph{Αλλά} στα ομοιογενή panel χωρίς Pinnacle ο
απλός μέσος όρος γίνεται ανταγωνιστικός, και στο opening ο VC\_BV μας κερδίζει σημαντικά. Η
τίμια διατύπωση είναι ότι η επιλεγμένη διαμόρφωση είναι η καλύτερη \emph{εν μέρει επειδή}
περιέχει τον καλύτερο odds-setter της αγοράς.}

\qa{Ξεπερνάς τον απλό μέσο όρο. Πόσο από αυτό είναι πραγματικά μάθηση;}
{\textbf{Λιγότερο απ' όσο φαίνεται, και το μέτρησα.} Υπό τη Sleeping Experts το διάνυσμα
βαρών κινείται ακόμη και με \emph{μηδενικό} βήμα, γιατί η προβολή επαναλαμβάνεται σε simplex
άλλης διάστασης κάθε φορά που αλλάζει το σύνολο των ενεργών experts --- στο $\eta=0$ τα βάρη
απέχουν ως 0,69 από το ομοιόμορφο. Βάζοντας ρητό control μηδενικού βήματος, το \textbf{41\%}
του περιθωρίου έναντι του μέσου όρου είναι αυτή η γεωμετρία και όχι ο αλγόριθμος.}
\trap{Αυτή είναι η ερώτηση που πρέπει να την πω εγώ πρώτος. Αν την πει η επιτροπή και δεν
την έχω προβλέψει, μοιάζει με κάτι που δεν είχα δει.}

\qa{Διάλεξες τη διαμόρφωση κοιτώντας τα αποτελέσματα. Δεν είναι κι αυτό επιλογή;}
{Ναι, και γι' αυτό έβαλα χρονολογική ουρά 25\% που δεν συμμετείχε στην επιλογή. Το
αποτέλεσμα είναι ότι στην ουρά και οι έξι διαμορφώσεις μαζεύονται σε λιγότερο από 0,00005. Η
ουρά δηλαδή \emph{επιβεβαιώνει την ισοπαλία, δεν βγάζει νικητή}, και το γράφω έτσι. Θα ήταν
ανέντιμο να παρουσιαστεί ως νίκη.}

\qa{Δοκίμασες πάρα πολλές ιδέες και οι περισσότερες απέτυχαν. Δεν είναι κι αυτό ψάρεμα;}
{Θα ήταν, αν ανέφερα μόνο όσες πέτυχαν. Αναφέρω και τις εννέα που απορρίφθηκαν, με τους
λόγους. Και κάθε υπερπαράμετρος επιλέγεται σε χρονολογικό πρόθεμα, ποτέ στην ουρά. Το
κριτήριο δεν είναι πόσες δοκίμασα αλλά αν η επιλογή έγινε με πληροφορία που δεν θα είχα τη
στιγμή της απόφασης.}

\qa{Οι συντελεστές βήματος των Hedge και FTRL δεν είναι πια βέλτιστοι. Γιατί δεν τους
ξαναρύθμισες;}
{Γιατί η επαναρύθμιση θα άλλαζε κάθε νούμερο στην εργασία και τα αποτελέσματα δεν θα ήταν
πλέον συγκρίσιμα μεταξύ τους. Το κόστος είναι ότι οι δύο αυτοί αλγόριθμοι αναφέρονται με
ελαφρώς \emph{υποβέλτιστο} βήμα, δηλαδή συντηρητικά --- δεν ευνοεί το συμπέρασμά μου, το
αντίθετο. Και είναι από μόνο του εύρημα: η προσθήκη δύο experts μετακινεί το βέλτιστο βήμα,
άρα ο συντονισμός είναι ιδιότητα της \emph{σύνθεσης του panel} και όχι του αλγορίθμου.}

\qa{Κάποιες εταιρείες δεν καλύπτουν ολόκληρη τη δεκαετία. Γιατί δεν τις πέταξες;}
{Γιατί η αναγωγή Sleeping Experts υπάρχει ακριβώς γι' αυτό: κάθε expert κρίνεται μόνο στους
γύρους όπου ήταν ενεργός, οπότε μια σειρά με μερική παρουσία δεν αραιώνει την εγγύηση των
υπολοίπων. Το να πετάξεις σειρές επειδή σταματούν νωρίς θα ήταν και ασύμμετρο, αφού
δεκατρείς σειρές του panel δεν υπάρχουν καθόλου στις πρώτες οκτώ σεζόν. Όπου η χρονική
παρουσία \emph{έχει} σημασία --- στα κατώφλια κάλυψης --- το κατώφλι ξαναϋπολογίζεται μέσα
στο παράθυρο που εξετάζεται.}

\qa{Ο έλεγχος ανάπτυξης έχει 8.151 αγώνες. Η «ισοπαλία» δεν είναι απλώς έλλειψη ισχύος;}
{Εν μέρει ναι, και το γράφω στους περιορισμούς. Τα διαστήματα εμπιστοσύνης είναι φαρδύτερα
από εκείνα της πλήρους σύγκρισης με 71.818 αγώνες, οπότε «ισοπαλία» σημαίνει εν μέρει
«λιγότερα δεδομένα». Αυτό που \emph{δεν} εξηγείται από έλλειψη ισχύος είναι ότι στο ίδιο
δείγμα οι τιμές ανοίγματος κερδίζονται σημαντικά. Η ίδια ισχύς αρκεί εκεί.}

% =====================================================================
\section{Οικονομικά}
% =====================================================================

\qa{Προβλέπεις καλύτερα. Γιατί δεν βγάζεις χρήματα;}
{Γιατί το κατώφλι δεν είναι «να είμαι καλύτερος από τον bookmaker» αλλά «να είμαι καλύτερος
\emph{κατά περισσότερο από το περιθώριό του}». Το στοίχημα απαιτεί $p\,o>1$ απέναντι στις
\emph{με γκανιότα} αποδόσεις. Η διαφορά ποιότητας που κέρδισα είναι μικρότερη από τη
γκανιότα, οπότε στις περισσότερες περιπτώσεις δεν ανοίγει θέση.}

\qa{Σε ποιους βγάζεις κέρδος, τελικά;}
{Σε έντεκα στόχους προκύπτουν \emph{δύο} σημαντικά κέρδη, στον WHC και στον IW, και μία
σημαντική ζημιά, στον PSC. Η ζημιά στον PSC είναι ο καλύτερος έλεγχος που έχω: είναι ο
καλύτερος προγνώστης του panel, άρα ένα κατώτερο μοντέλο \emph{οφείλει} να χάνει απέναντί
του. Ο έλεγχος λειτουργεί και προς τις δύο κατευθύνσεις.}

\qa{Άρα υπάρχει εκμεταλλεύσιμη άκρη;}
{Στατιστικά ναι, πρακτικά όχι απαραίτητα. Οι δύο εταιρείες που κερδίζονται είναι ακριβώς οι
«μαλακές», αυτές που απευθύνονται σε ερασιτέχνες --- και αυτές είναι οι πρώτες που κλείνουν
τους κερδισμένους λογαριασμούς. Η προσομοίωση δεν μοντελοποιεί όρια λογαριασμού. Η πρόσβαση
και η εκμεταλλεύσιμη άκρη είναι αντιστρόφως ανάλογες από την ίδια τη δομή του κλάδου.}

\qa{Γιατί δεν ξεπερνιέται η Pinnacle;}
{Όχι επειδή έχει καλύτερο μοντέλο, αλλά επειδή έχει διαφορετικό \emph{επιχειρηματικό
μοντέλο}: χαμηλό περιθώριο, υψηλός τζίρος, και δεν περιορίζει τους κερδισμένους. Αυτό
σημαίνει ότι η τιμή της δεν είναι η γνώμη μιας εταιρείας αλλά συσσωρευμένη τιμή από
πληροφορημένο χρήμα. Είναι ισχυρισμός για τον κλάδο, όχι αποτέλεσμα των δεδομένων μου, και
το λέω έτσι.}

% =====================================================================
\section{Η νέα δουλειά}
% =====================================================================

\qa{Τι είναι το bandit πείραμα και τι έδειξε;}
{Ρωτά τι απομένει αν ο αλγόριθμος βλέπει μόνο \emph{έναν} αριθμό τον γύρο --- την απώλεια
της δικής του πρόβλεψης --- αντί για ολόκληρο το διάνυσμα απωλειών. Έδειξε δύο πράγματα. Ο
εκτιμητής ενός σημείου \textbf{δεν μαθαίνει τίποτα} εδώ, και ξέρω γιατί: η ευθυγράμμισή του
με την αληθινή κλίση είναι 0,002, οπότε ο ωφέλιμος λόγος σε 76.584 γύρους βγαίνει 0,55,
κάτω από τη μονάδα. Και το δεύτερο, σημαντικότερο, είναι το control μηδενικού βήματος που
αποκάλυψε το 41\% της προβολής.}

\qa{Δοκίμασες να το διορθώσεις;}
{Ναι, με mini-batching, που είναι η μία παρέμβαση που δικαιολογεί το επιχείρημα της
διασποράς. Σάρωσα τέσσερα μεγέθη batch επί πέντε τάξεις μεγέθους του βήματος. Κανένας
συνδυασμός δεν δίνει εσωτερικό βέλτιστο. Το batching κάνει ακριβώς ό,τι λέει η θεωρία,
ανέχεται μεγαλύτερο βήμα χωρίς ζημιά, αλλά δεν παράγει κέρδος γιατί το σήμα δεν υπάρχει.}

\qa{Γιατί η μάθηση του πολλαπλασιαστή Kelly δουλεύει ενώ η μάθηση του πονταρίσματος όχι;}
{Είναι θέμα \emph{comparator}. Ένα φράγμα regret για μαθητή πονταρίσματος υπόσχεται απέναντι
στο καλύτερο \emph{σταθερό} ποντάρισμα. Ο κλειστός τύπος όμως δεν ανήκει σε αυτή την κλάση:
παράγει διαφορετικό ποντάρισμα κάθε γύρο από το $\hat{p}$ και την απόδοση, που είναι γνωστά
\emph{πριν} το στοίχημα. Ο μαθητής επιπλέον πετά το $\hat{p}$ που τα δύο προηγούμενα στάδια
μόλις παρήγαγαν. Μαθαίνοντας τον \emph{πολλαπλασιαστή}, η πληροφορία του γύρου μένει μέσα
και ο comparator «καλύτερο σταθερό $\lambda$» είναι ακριβώς το $1/4$ που καρφώνουμε.}

\qa{Και τι βρήκες;}
{Ισοπαλία με το σταθερό $1/4$, 17 ισοπαλίες και 2 νίκες σε 20 κελιά, αλλά μπροστά στο
κεφάλαιο όπου υπάρχει πραγματική άκρη. Το ενδιαφέρον είναι \emph{τι μαθαίνει}: το $\lambda$
βγαίνει 0,37 στον ζημιογόνο PSC και 0,79 στον κερδοφόρο WHC. Ανεβάζει το ρίσκο όπου υπάρχει
περιθώριο και το μαζεύει όπου δεν υπάρχει, χωρίς να του πω ποιος είναι ο αντίπαλος. Και είναι
το μόνο σημείο όλης της εργασίας όπου η θεωρητική σταθερά βγήκε σωστή χωρίς διόρθωση.}

\qa{Γιατί απέτυχε η εναλλαγή καθεστώτων;}
{Με διδακτικό τρόπο. Στην πρώτη εκδοχή η μέση πεποίθηση έμεινε \emph{ακριβώς} 0,500 και στα
δύο καθεστώτα: δεν διαφοροποιήθηκαν ποτέ, γιατί βλέπουν το ίδιο διάνυσμα απωλειών και από
συμμετρική αφετηρία ακολουθούν ίδιες τροχιές. Ήταν δηλαδή απλώς OGD με μισό βήμα. Μετά τη
διόρθωση ισοπαλεί, αλλά επιλέγεται $\tau=0$ και η τυπική απόκλιση της πεποίθησης είναι
μηδέν --- άρα η αλυσίδα αποδεδειγμένα δεν λειτουργεί, και ό,τι μένει είναι ensembling.}

% =====================================================================
\section{Στατιστικά}
% =====================================================================

\qa{Γιατί moving block bootstrap και όχι $t$-test;}
{Γιατί οι διαδοχικές απώλειες δεν είναι ανεξάρτητες: αγώνες της ίδιας αγωνιστικής και
περιόδου μοιράζονται δυσκολία. Σε θετικά συσχετισμένη ακολουθία ο τύπος της ανεξαρτησίας
\emph{υποτιμά} τη διασπορά, τα διαστήματα βγαίνουν πολύ στενά και εμφανίζονται ψευδώς
σημαντικά αποτελέσματα. Το bootstrap μπλοκ δειγματοληπτεί διαδοχικές παρατηρήσεις μαζί,
οπότε διατηρεί την τοπική εξάρτηση.}

\qa{Πώς διάλεξες το μήκος μπλοκ;}
{Είναι ανταλλαγή μεροληψίας--διασποράς με αποτέλεσμα $b\propto n^{1/3}$, που για το δικό μου
$n$ δίνει περίπου 42 αγώνες, δηλαδή μιάμιση αγωνιστική. Αξίζει να σημειωθεί ότι το ίδιο
μέγεθος εμφανίστηκε αυτόνομα ως ορίζοντας εξάρτησης στη μελέτη του βήματος, από εντελώς
διαφορετικό συλλογισμό.}

\qa{Τι σημαίνει το πρόσημο στους πίνακές σου;}
{Προσοχή, αλλάζει ανά πίνακα. Στο log-loss ορίζω $d=\ell(A)-\ell(B)$, οπότε \emph{αρνητικό}
σημαίνει ότι ο πρώτος είναι καλύτερος. Στους πίνακες λογαριθμικής \emph{ανάπτυξης} κεφαλαίου
ισχύει το αντίθετο, θετικό σημαίνει καλύτερος. Το έχω σημειώσει σε κάθε πίνακα.}

% =====================================================================
\section{Τα τρία συμπεράσματα, αν με ρωτήσουν «τι έμαθες»}
% =====================================================================

\qa{Ποιο είναι το πιο χρήσιμο συμπέρασμα για κάποιον που θα εφάρμοζε το πλαίσιο αλλού;}
{Ότι η \emph{ανομοιογένεια} του panel είναι προϋπόθεση, όχι λεπτομέρεια. Η μαθημένη στάθμιση
αποδίδει μόνο όταν οι experts διαφέρουν ουσιωδώς σε ποιότητα. Σε ομοιογενές σύνολο ο απλός
μέσος όρος είναι τουλάχιστον εξίσου καλός, γιατί η προσαρμογή πληρώνει θόρυβο χωρίς
αντίστοιχο όφελος. Το κριτήριο για καλό σύνολο experts δεν είναι το πλήθος τους αλλά αν
\emph{παράγονται από διαφορετικές διαδικασίες}.}

\qa{Ποιο θα κρατούσες ως το πιο ειλικρινές;}
{Ότι ένα σημαντικό μέρος της φαινομενικής υπεροχής δεν ήταν μάθηση. Το βρήκα επειδή έβαλα
control που δεν χρειαζόταν να βάλω, και άλλαξε συμπέρασμα που είχα ήδη γράψει.}

\qa{Και το πιο πρακτικό;}
{Ότι καλύτερη πρόβλεψη δεν σημαίνει κέρδος. Η διάκριση ανάμεσα στο να είσαι καλύτερος
προγνώστης και στο να μπορείς να το εκμεταλλευτείς είναι το ουσιαστικό συμπέρασμα, και
ισχύει πολύ πέρα από το στοίχημα.}

\vspace{1.5em}
\hrule
\vspace{0.8em}
{\small
\textbf{Τρεις κανόνες για την παρουσίαση.} Πρώτον, κάθε αδυναμία να ειπωθεί από εμένα πριν
τη βρει η επιτροπή --- ιδίως το 41\% της προβολής και το ότι η Pinnacle είναι μέσα στο
μείγμα. Δεύτερον, να μη λέγεται «ξεπερνάμε την αγορά»: η σωστή διατύπωση είναι ότι
\emph{φτάνουμε} την αιχμή των bookmakers. Τρίτον, για κάθε ισοπαλία να λέγεται αν οφείλεται
σε ίδια ποιότητα ή σε λίγα δεδομένα· συνήθως είναι και τα δύο.
}

\end{document}
